
%%%%%%%%%%%%%%%%%%%%%%%%%%%%%%%%%%%%%%%%%%%%%%%%%%%%%%%%%%%%%%%%%%%%%%%%%%%%%%%%%%
%\subsection{ Gauge  transformations}
%\label{sec:velocity_gauge}
%
%In this section we discuss traansformation og Time-Dependent Kohn-Sham  equation  under differet gauge representation within Coulomb gauge condtion.
%
%For finite systems, the relevant observable is the
%time-dependent dipole moment $\vb d(t) = \int \vb r\,
%\rho(\vb r,t)\, d\vb r$, from which the polarizability
%and absorption cross section follow by Fourier transform.
%For periodic systems, however, the scalar-potential form
%$\Phi = -\vb E(t)\cdot\vb r$ is not compatible with
%periodic boundary conditions, because the position
%operator $\vb r$ is discontinuous at the cell boundary.
%This motivates the use of the velocity-gauge (VG) formulation.
%In the velocity gauge, the position operator no longer
%appears in the Hamiltonian and all terms are compatible
%with periodic boundary conditions.
%Correspondingly, the physical observable used to extract
%the optical response is the time-dependent current density
%$\vb j(\vb r,t)$ rather than the dipole moment, since
%the current is well defined under periodic boundary
%conditions whereas $\vb d(t)$ is not.


% ======================================================================
%  Gauge Transformation from Length to Velocity Gauge
%  Follows sec_em_interaction_blue_review-v3.tex and uses its labels:
%    sec:em-interaction, eq:em-tdks, eq:em_E_from_A+Phi, eq:em_B_from_A,
%    eq:gauge-rules-A (NEW: add to the A-line of the gauge-rules align),
%    eq:gauge-rules, eq:EM_dipolar.
%  Defines \label{sec:velocity_gauge}, as required by that file.
%  Remaining TODO: label of the current-operator section (sec:current).
% ======================================================================

\subsection{Gauge Transformation from Length to Velocity Gauge}
\label{sec:velocity_gauge}

Within the electric-dipole (long-wavelength) approximation
a spatially uniform external  electric field $\vb{E}_\mathrm{ext}(t)$ can be represented by two
gauge-equivalent pairs of potentials. In the length gauge (LG) the field enters
through the linear in position dipolar scalar potential of Eq.~\eqref{eq:EM_dipolar}
\begin{equation}
\vb{A}_L(\vb{r},t) = 0,  \qquad
	\Phi_L(\vb{r},t) = -\,\vb{E}_\mathrm{ext}(t)\cdot\vb{r}
	  \label{eq:LG_potentials}
\end{equation}
whereas in the velocity gauge (VG) it enters through a spatially uniform,
time-dependent vector potential,
\begin{equation}
  \vb{A}_V(\vb{r},t) = \vb{A}(t), \qquad
  \Phi_V(\vb{r},t) = 0, \qquad
  \label{eq:VG_potentials}
\end{equation}
where the relation between vector potential   $\vb{A}(t)$  and  external electric field $\vb{E}_\mathrm{ext}(t)$  
is given by
\begin{equation}
  \vb{E}_\mathrm{ext}(t) = -\frac{1}{c}\frac{d\vb{A}(t)}{dt}.
\end{equation}
Inserting either pair, Eqs.~\eqref{eq:LG_potentials} or \eqref{eq:VG_potentials},
into Eqs.~\eqref{eq:em_E_from_A+Phi} and \eqref{eq:em_B_from_A}
gives the same uniform $\vb{E}=\vb{E}_\mathrm{ext}(t)$ and $\vb{B}=0$. The LG form is natural
for finite systems, but $\Phi_L$ is not lattice periodic, whereas the VG
Hamiltonian contains only the uniform $\vb{A}(t)$ and is compatible with
Born--von K\'arm\'an boundary conditions. 
%In LG,  the   optical response is usually  extracted  from   the time dependent  dipole moment  $\vb d(t)$.
In VG,  the   optical response is extracted  from   the 
  time-dependent current density
$\vb j(\vb r,t)$ which is, contrary to  the dipole moment  $\vb d(t)$ used in LG, 
 well defined under periodic boundary conditions.
The two representations are connected
by the G\"oppert-Mayer gauge transformation.\cite{GoeppertMayer-1931,Mattiat-Luber-2022}
Below we directly obtain gauge function for transformation from LG to VG and discuss the transformation of
time-dependent Kohn--Sham (TDKS) equations, corresponding operators, and pseudopotentials.



\paragraph{Gauge function.}
The gauge function for LG to VG transformation can be obtained directly from the potential transformation rules,
Eqs.~\eqref{eq:gauge-rules-A} and \eqref{eq:gauge-rules}, and a requirement that the transformed scalar potential $\Phi_V$  vanishes
\begin{equation}
  \Phi_V = \Phi_L - \frac{1}{c}\frac{\partial\chi}{\partial t} = 0
 % \quad\Longrightarrow\quad
 % \frac{\partial\chi(\vb{r},t)}{\partial t} = -c\,\vb{r}\cdot\vb{E}(t).
  \label{eq:chi_condition}
\end{equation}
Inserting $\Phi_L$  from  Eq.~\eqref{eq:EM_dipolar} into the last expression and solving  for  
$\chi(\vb r,t)$  gives 
\begin{equation}
  \chi(\vb{r},t) = \vb{r}\cdot\vb{A}(t) + \chi_0(\vb{r}),
  \qquad
	\vb{A}(t) = -c\int_{t_0}^{t}\vb{E}_\mathrm{ext}(t')\,dt',
  \label{eq:chi_solution}
\end{equation}
%where $\chi_0(\vb r)$ is an arbitrary time-independent function. Choosing
%$\chi_0 = 0$ gives $\vb A' = \vb A(t)$, i.e., the VG potentials of
%Eq.~\eqref{eq:VG_potentials} with $\vb A(t_0) = 0$. Since
%$\chi = \vb r\cdot\vb A(t)$ is linear in $\vb r$, $\nabla^2\chi = 0$ and both
%gauges satisfy the Coulomb condition, $\nabla\cdot\vb A = 0$; and since
%$\nabla\times\nabla\chi = 0$, $\vb B = \nabla\times\vb A$ is unchanged and
%vanishes in both gauges. We note that static gauge transformations are also
%useful in practice, for example to remove the spurious gauge-origin dependence
%of finite-basis calculations of magnetic susceptibilities and NMR
%shieldings.\cite{Ditchfield-1974,Keith-Bader-1993,pickard-mauri-2001}
%
where $\chi_0(\vb{r})$ is an arbitrary time-independent function. The
transformed vector potential is
$\vb{A}' = \vb{A}_L + \nabla\chi = \vb{A}(t) + \nabla\chi_0$, and the
choice $\chi_0 = 0$ yields exactly the VG potentials of
Eq.~\eqref{eq:VG_potentials}, with $\vb{A}(t_0)=0$. 
%Since $\nabla^2\chi = 0$ and $\nabla\times\nabla\chi = 0$, both representations satisfy
%the Coulomb-gauge condition and neither carries a magnetic field.
With $\chi_0 = 0$, the gauge function  $\chi(\vb r,t)$ in Eq.~\eqref{eq:chi_solution} is  linear in $\vb r$, so $\nabla^2\chi = 0$ and both
representations satisfy the Coulomb-gauge condition, $\nabla\cdot\vb A = 0$.
Since $\nabla\times\nabla\chi = 0$, the gauge transformation leaves
$\vb B = \nabla\times\vb A$ unchanged, and the magnetic field vanishes in both
gauges.
We note that static gauge transformations, where $\chi_0(\vb r)\neq 0$,  are also useful in practice, for
example to remove the spurious gauge-origin dependence of finite-basis
calculations of magnetic susceptibilities and NMR
shieldings.\cite{Ditchfield-1974,Keith-Bader-1993,pickard-mauri-2001}

The spatial uniformity of $\vb{A}(t)$ is a direct consequence of the
\emph{linear} position dependence of $\Phi_L$. The same procedure can be applied
formally to any scalar potential, $\chi = c\int_{t_0}^{t}\Phi(\vb{r},t')\,dt'$,
but if $\Phi$ contains higher-order polynomial terms in $\vb{r}$ (for example
a quadrupolar field), $\vb{A}'=\nabla\chi$ is itself position dependent. That is,
the position operator is removed from the scalar potential only to reappear in the
vector potential, and no lattice-periodic representation is obtained.





%\paragraph{Length- to velocity-gauge TDKS Hamiltonian.}

\paragraph{Gauge invariance of observables.}
%
%The gauge function of Eq.~\eqref{eq:chi_solution} determines how the vector and scalar potentials
%change, but not yet how the Kohn--Sham orbitals and the operators acting on them
%must change, so that the TDKS equations describe the same physics in the new
%gauge. We derive these rules for a general $\chi(\vb r,t)$ from a single
%requirement: a gauge transformation must not alter any physically measurable
%quantity. If $\hat O$ represents such a quantity and $\hat O'$ is its
%representation in the transformed gauge, then for every Kohn--Sham orbital
%
%
The gauge function of Eq.~\eqref{eq:chi_solution} determines how the vector and scalar potentials
change, but not how the Kohn--Sham orbitals and the operators transform.
Here we disscuss these rules deriving them for a general $\chi(\vb r,t)$ from 
a requirement that a gauge transformation must not alter any physically measurable
quantity. 
%If $\hat O$ represents such a quantity and $\hat O'$ is its
%representation in the transformed gauge, then for every Kohn--Sham orbital
%It is required  that transformation must not alter any physically measurable quantity. 
In particular  if operator $\hat O$ represents such a quantity and $\hat O'$ is its representation 
in the transformed gauge, then for every Kohn–Sham orbital
\begin{equation}
%  \hat O'=\hat U_\chi\hat O\hat U_\chi^\dagger,
%  \qquad
  \langle\Psi'_j|\hat O'|\Psi'_j\rangle
  =
  \langle\Psi_j|\hat O|\Psi_j\rangle.
  \label{eq:observable_invariance}
\end{equation}
%\myblue{This requirement fixes, in turn, the transformation of the orbitals, the phase
%constant, and the TDKS Hamiltonian.}
Taking, for example, the density operator $\hat { n}(\vb r_0) = \delta(\vb{\hat r}-\vb r_0)$ ,
which has the same form in every gauge, Eq.~\eqref{eq:observable_invariance} 
%where, for example, $\hat O$ operator is a density operator $\hat n(\vb r_0) = \delta(\vb{\hat r}-\vb r_0)$.
%Equation~\eqref{eq:observable_invariance} therefore
requires $|\Psi_j'(\vb r,t)|^2 = |\Psi_j(\vb r,t)|^2$ at every point, which is
possible only if the orbitals differ by a local phase,
\begin{equation}
  \Psi_j'(\vb r,t) = e^{\imath\theta(\vb r,t)}\,\Psi_j(\vb r,t).
  \label{eq:local_phase}
\end{equation}
The gauge transformation is a one-particle operation that does not depend on
which states are occupied, so the same phase applies to every orbital. The
phase must vanish when $\chi = 0$. % and must depend locally on $\chi$. 
Moreover, two successive transformations with $\chi_1$ and $\chi_2$ are equivalent to a
single transformation with $\chi_1+\chi_2$, so the phases must add. These
conditions require $\theta$ to be linear in $\chi$:
\begin{equation}
  \Psi_j' = \hat U_\chi\,\Psi_j, \qquad
  \hat U_\chi(t) = \exp\!\big[\imath\beta\,\chi(\vb r,t)\big],
  \label{eq:U_beta}
\end{equation}
with a real constant $\beta$ still to be determined. Because $\chi$ is real,
$\hat U_\chi$ is unitary, $\hat U_\chi^\dagger\hat U_\chi = \hat I$.
Equation~\eqref{eq:observable_invariance} is then satisfied for arbitrary
states if observables transform as
\begin{equation}
  \hat O' = \hat U_\chi\,\hat O\,\hat U_\chi^\dagger .
  \label{eq:observable_rule}
\end{equation}

%%%%%%%%%%%%%%%%%%%%

\paragraph{Canonical and kinetic momentum.}
Here we show that the constant $\beta$ in Eq.~\eqref{eq:U_beta} is fixed by the
momentum. The canonical momentum $\vb{\hat p} = -\imath\hbar\nabla$ is
conjugate to position, $[\hat r_\alpha,\hat p_\beta] = \imath\hbar\,\delta_{\alpha\beta}$,
and keeps this form in every gauge. Its expectation value, however, changes
under the orbital transformation of Eq.~\eqref{eq:U_beta}. That is 
$\hat U_\chi^\dagger\,\vb{\hat p}\,\hat U_\chi = \vb{\hat p} + \hbar\beta\nabla\chi$,
and
\begin{equation}
  \langle\Psi_j'|\vb{\hat p}|\Psi_j'\rangle
  = \langle\Psi_j|\vb{\hat p}|\Psi_j\rangle
  + \hbar\beta\,\langle\Psi_j|\nabla\chi|\Psi_j\rangle .
  \label{eq:p_shift}
\end{equation}
As shown below  $\beta\neq 0$, so  the canonical momentum is gauge dependent and
is not an observable.

The kinetic (mechanical) momentum,
\begin{equation}
  \hat{\boldsymbol{\pi}} = \vb{\hat p} - \frac{q_e}{c}\,\vb A ,
  \label{eq:kinetic_momentum}
\end{equation}
is related to the velocity of the electron, which is a physical observable. For
a Hamiltonian containing only local potentials, the Heisenberg velocity
operator is
\begin{equation}
  \vb{\hat v} = \frac{\imath}{\hbar}\,[\hat H,\vb{\hat r}]
             = \frac{\hat{\boldsymbol{\pi}}}{m_e} ;
  \label{eq:velocity_operator}
\end{equation}
the nonlocal pseudopotential adds a commutator contribution, discussed in
Sec.~\ref{sec:current}. Because the velocity, and hence the current, is
measurable, $\hat{\boldsymbol{\pi}}$ must satisfy
Eq.~\eqref{eq:observable_invariance}, with
$\hat{\boldsymbol{\pi}}' = \vb{\hat p} - (q_e/c)\vb A'$ in the transformed gauge.
Using $\vb A' = \vb A + \nabla\chi$ and Eq.~\eqref{eq:p_shift}, we obtain
\begin{equation}
  \langle\Psi_j'|\hat{\boldsymbol{\pi}}'|\Psi_j'\rangle
  = \langle\Psi_j|\hat{\boldsymbol{\pi}}|\Psi_j\rangle
  + \Big(\hbar\beta - \frac{q_e}{c}\Big)\langle\Psi_j|\nabla\chi|\Psi_j\rangle .
  \label{eq:pi_expectation}
\end{equation}
Since $\chi$ is arbitrary, Eq.~\eqref{eq:observable_invariance}
requires that the term in parentheses vanishes, which yields
$\beta = q_e/(\hbar c)$ and
\begin{equation}
  \hat U_\chi(t) = \exp\!\Big[\frac{\imath q_e}{\hbar c}\,\chi(\vb r,t)\Big] .
  \label{eq:U_final}
\end{equation}
With this value of $\beta$, the canonical and kinetic momentum operators
transform as
\begin{equation}
  \hat U_\chi\,\vb{\hat p}\,\hat U_\chi^\dagger = \vb{\hat p} - \frac{q_e}{c}\nabla\chi,
  \qquad
  \hat U_\chi\,\hat{\boldsymbol{\pi}}[\vb A]\,\hat U_\chi^\dagger
  = \hat{\boldsymbol{\pi}}[\vb A'] .
  \label{eq:momentum_transform}
\end{equation}
The shift of canonical momentum in Eq.~\eqref{eq:p_shift} 
reflects the fact that the vector potential contained in
%thus becomes $(q_e/c)\langle\nabla\chi\rangle$,
%reflecting the vector potential contained in
$\vb{\hat p} = \hat{\boldsymbol{\pi}} + (q_e/c)\vb A$  is defined only up
to $\nabla\chi$. The kinetic momentum, in contrast, transforms covariantly, 
keeping its form with the transformed potential, consistent with
Eq.~\eqref{eq:observable_rule}.






%%%%%%%%%%%%%%%%%%%
\paragraph{Time-dependent Kohn--Sham equation.}
Starting from the time-dependent Kohn--Sham equation,
\begin{equation}
  \imath\hbar\,\partial_t\Psi_j = \hat H(t)\,\Psi_j ,
  \label{eq:tdks_general}
\end{equation}
we apply the product rule to the transformed orbital:
\begin{equation}
\begin{split}
  \imath\hbar\,\partial_t\Psi_j'
  &= \imath\hbar\,(\partial_t\hat U_\chi)\,\Psi_j
   + \hat U_\chi\,\imath\hbar\,\partial_t\Psi_j \\
  &= \Big[\imath\hbar\,(\partial_t\hat U_\chi)\hat U_\chi^\dagger
   + \hat U_\chi\hat H\hat U_\chi^\dagger\Big]\Psi_j' .
\end{split}
\label{eq:tdks_product_rule}
\end{equation}
%Evaluating the time derivative of $\hat U_\chi$, Eq.~\eqref{eq:U_final},  gives
%$  \imath\hbar\,(\partial_t\hat U_\chi)\hat U_\chi^\dagger
%  = -\frac{q_e}{c}\,\partial_t\chi$.
Thus the transformed orbitals obey $\imath\hbar\,\partial_t\Psi_j' = \hat H'\Psi_j'$, with
\begin{equation}
\begin{split}
\hat H' = \hat U_\chi\hat H\hat U_\chi^\dagger
        + \imath\hbar\,(\partial_t\hat U_\chi)\hat U_\chi^\dagger  \\
        = \hat U_\chi\hat H\hat U_\chi^\dagger - \frac{q_e}{c}\,\partial_t\chi .
\end{split}
	\label{eq:H_transformed}
\end{equation}

Because $\hat U_\chi$ depends on time, the Hamiltonian does not transform
according to the observable rule, Eq.~\eqref{eq:observable_rule}. For
$\partial_t\chi\neq 0$, the expectation value of $\hat H$ is gauge dependent:
$\langle\Psi_j'|\hat H'|\Psi_j'\rangle = \langle\Psi_j|\hat H|\Psi_j\rangle
- (q_e/c)\langle\Psi_j|\partial_t\chi|\Psi_j\rangle$.
%\myblue{and only the mechanical
%energy, $\hat\varepsilon = \hat{\boldsymbol{\pi}}^2/2m_e + \hat V_{\mathrm{KS}}$,
%transforms covariantly and is a gauge-invariant observable.
%}

%%%%%%%%%%%%%%%%%%%%%%%%%%%%%%%%%%%%%%%%%%%`
\paragraph{Hamiltonian transformation.}
The individual terms of $\hat H$ transform as follows. The external scalar
potential and the local Kohn--Sham potentials (namely the local ionic, Hartree,
and exchange--correlation terms)  are multiplicative operators, so they commute
with $\hat U_\chi$. Moreover, the Hartree and exchange--correlation potentials
depend only on the density, which is gauge invariant by
Eq.~\eqref{eq:observable_invariance}, so they are the same functionals in both
gauges. Combined with the time-derivative term of Eq.~\eqref{eq:H_transformed},
the scalar-potential contribution becomes
$q_e\Phi - (q_e/c)\,\partial_t\chi = q_e\Phi'$, in agreement with
Eq.~\eqref{eq:gauge-rules}. The kinetic term transforms through the covariance
of the kinetic momentum, Eq.~\eqref{eq:momentum_transform}:
\begin{equation}
  \hat U_\chi\,\hat{\boldsymbol{\pi}}^2\,\hat U_\chi^\dagger
  = \big(\hat U_\chi\hat{\boldsymbol{\pi}}\,\hat U_\chi^\dagger\big)
    \big(\hat U_\chi\hat{\boldsymbol{\pi}}\,\hat U_\chi^\dagger\big)
  = \hat{\boldsymbol{\pi}}'^{\,2},
  \qquad
  \hat{\boldsymbol{\pi}}' = \vb{\hat p} - \frac{q_e}{c}\vb A' .
  \nonumber
\end{equation}

The transformed Hamiltonian therefore has the same form as the original one,
with $\vb A$ and $\Phi$ replaced by $\vb A'$ and $\Phi'$:
\begin{equation}
  \hat H' = \frac{1}{2m_e}\Big[\vb{\hat p} - \frac{q_e}{c}\vb A'\Big]^2
          + q_e\Phi' + V_{\mathrm{loc}}[\rho]
          + \hat U_\chi\hat V_{\mathrm{NL}}\hat U_\chi^\dagger, 
  \label{eq:H_general_transformed}
\end{equation}
i.e., the TDKS equation is form invariant.
The only exception is the nonlocal pseudopotential. It does not commute with
$\hat U_\chi$ and must be transformed explicitly; the rules for multiplicative,
differential, and nonlocal operators are summarized in
Appendix~\ref{app:gauge_rules}, Table~\ref{tab:gauge_operator_classes}.
%Appendix~\ref{sec:gauge_rules}, Table~\ref{tab:gauge_operator_classes}.

%Form invariance does not contradict the gauge dependence of $\langle\hat H\rangle$
%found above: $\hat H'$ has the same form as $\hat H$, but contains the
%transformed scalar potential $\Phi'\neq\Phi$. Only the mechanical energy,
%$\hat\varepsilon = \hat{\boldsymbol{\pi}}^2/2m_e + \hat V_{\mathrm{KS}}$, which
%excludes the external scalar-potential term, transforms covariantly and is
%gauge invariant.
%
%
%Because of the $\partial_t\chi$ term in Eq.~\eqref{eq:H_transformed}, the
%expectation value of $\hat H$ itself is gauge dependent,
%$\langle\Psi_j'|\hat H'|\Psi_j'\rangle = \langle\Psi_j|\hat H|\Psi_j\rangle
%- (q_e/c)\langle\Psi_j|\partial_t\chi|\Psi_j\rangle$. Only the mechanical
%energy, $\hat\varepsilon = \hat{\boldsymbol{\pi}}^2/2m_e + \hat V_{\mathrm{KS}}$,
%which excludes the external scalar-potential term, transforms covariantly and
%is gauge invariant.


%%%%%%%%%%%%%%%%%%%%%%%%%%%%%%%%%%%%%%%%%%%%%%%%%%%%
\paragraph{Nonlocal pseudopotential.}
Unlike the local potentials, the nonlocal pseudopotential is an integral operator,
$\hat V_{\mathrm{NL}}\Psi(\vb r) = \int V_{\mathrm{NL}}(\vb r,\vb r')\,\Psi(\vb r')\,d^3r'$
%$(\hat V_{\mathrm{NL}}\Psi)(\vb r) = \int V_{\mathrm{NL}}(\vb r,\vb r')\,\Psi(\vb r')\,d^3r'$,
whose kernel connects different points. 
The two gauge phases in $\hat U_\chi\hat V_{\mathrm{NL}}\hat U_\chi^\dagger$ are
therefore evaluated at different arguments and do not cancel.
For a  general  gauge function  the transformed  pseudopotential becomes
\begin{equation}
  V_{\mathrm{NL}}(\vb r,\vb r';t)
	= \exp\!\left[\frac{\imath q_e}{\hbar c}\, \big(\chi(\vb r) -\chi(\vb r') \big)\right]
    V_{\mathrm{NL}}(\vb r,\vb r') .
  \label{eq:VNL_general}
\end{equation}
In VG, with $\chi = \vb r\cdot\vb A(t)$, the transformed  non-local pseudpotential becomes
\begin{equation}
  V^{V}_{\mathrm{NL}}(\vb r,\vb r';t)
  = \exp\!\left[\frac{\imath q_e}{\hbar c}\,\vb A(t)\cdot(\vb r-\vb r')\right]
    V_{\mathrm{NL}}(\vb r,\vb r') .
  \label{eq:VNL_kernel}
\end{equation}
which  to first order in $\vb A$ is
\begin{equation}
  \hat V^{V}_{\mathrm{NL}}(t)
  = \hat V_{\mathrm{NL}}
  + \frac{\imath q_e}{\hbar c}\,\vb A(t)\cdot[\vb{\hat r},\hat V_{\mathrm{NL}}]
  + \mathcal{O}(A^2).
  \label{eq:VNL_linear}
\end{equation}
Because in VG the phase depends only on $\vb r-\vb r'$, the transformed operator
remains invariant under simultaneous lattice translation of both arguments and
retains the periodicity of the crystal.
\myblue{ [WENCHANG or JERRY, check this]
For the separable Kleinman--Bylander form,
$\hat V_{\mathrm{NL}} = \sum_{I,lm}|\beta_{I,lm}\rangle D_{I,lm}\langle\beta_{I,lm}|$,
Eq.~\eqref{eq:VNL_kernel} amounts to multiplying each projector by a gauge phase,
$\beta_{I,lm}(\vb r) \to e^{\imath q_e\vb A(t)\cdot\vb r/\hbar c}\,\beta_{I,lm}(\vb r)$. }
%This transformation is required for gauge consistency: transforming the kinetic
%term while leaving $\hat V_{\mathrm{NL}}$ unchanged breaks gauge
%invariance.\cite{IsmailBeigi-2001-EM,pickard-mauri-2001,Mattiat-Luber-2022}
The commutator term $[\vb{\hat r},\hat V_{\mathrm{NL}}] $  in Eq.~\eqref{eq:VNL_linear} also  enters the
nonlocal contribution to the velocity and current operators (see discussion  for Eq.~\eqref{eq:velocity_operator}
and Sec.~\ref{sec:current}). 
%Its treatment in the RMG perturbation is described in Sec.~\ref{sec:impl_perturbation}.

%%%%%%%%%%%%%%%%%
\paragraph{Velocity-gauge Hamiltonian and density matrix.}
Inserting the VG potentials, Eq.~\eqref{eq:VG_potentials}, and the transformed
nonlocal pseudopotential, Eq.~\eqref{eq:VNL_kernel}, into
Eq.~\eqref{eq:H_general_transformed} gives the TDKS Hamiltonian in the velocity
gauge,
\begin{equation}
  \hat H_V(t) = \frac{1}{2m_e}\Big[\vb{\hat p} - \frac{q_e}{c}\vb A(t)\Big]^2
              + V_{\mathrm{loc}}[\rho](\vb r,t) + \hat V^{V}_{\mathrm{NL}}(t) .
  \label{eq:H_VG}
\end{equation}
The dipolar term of the LG Hamiltonian, $q_e\Phi_L = -q_e\vb E_{\mathrm{ext}}(t)\cdot\vb r$,
is absent because $\chi$ was chosen to make $\Phi' = 0$, Eq.~\eqref{eq:chi_condition}:
the time-derivative term of Eq.~\eqref{eq:H_transformed},
$-(q_e/c)\,\partial_t\chi = q_e\vb E_{\mathrm{ext}}(t)\cdot\vb r$, cancels it
exactly. For an electron, $q_e = -e$, and the kinetic momentum is
$\vb{\hat p} + e\vb A(t)/c$.

 The one-particle density matrix,
$\hat P = \sum_j f_j|\Psi_j\rangle\langle\Psi_j|$, transforms as
$\hat P' = \hat U_\chi\hat P\,\hat U_\chi^\dagger$; its kernel acquires the
two-point phase of Eq.~\eqref{eq:VNL_general}, so only its diagonal, the
density $\rho(\vb r) = P(\vb r,\vb r)$, is gauge invariant, while traces with
consistently transformed observables are unchanged. The Liouville--von Neumann
equation retains its form,
$\imath\hbar\,\partial_t\hat P' = [\hat H',\hat P']$, with $\hat H'$ of
Eq.~\eqref{eq:H_transformed}. In practice, no such transformation is required:
because $\vb A(t_0) = 0$ implies $\chi(t_0) = 0$ and $\hat U_\chi(t_0) = 1$, the
field-free ground state and its density matrix are identical in both gauges. A
calculation that starts from the ground state and introduces the field only
through $\vb A(t)$ is therefore carried out in the velocity gauge from the
outset, and the ground-state density matrix serves directly as its initial
condition.


%%%%%%%%%%%%%%%%%%%%%%%%%%%%%%%%%%%%%%%%%%%%
\paragraph{Periodic representation CHECK THIS.}
Since every term in Eq.~\eqref{eq:H_VG} is lattice periodic, the VG orbitals
can be written in Bloch form,
$\Psi_{n\vb k}(\vb r,t) = e^{\imath\vb k\cdot\vb r}u_{n\vb k}(\vb r,t)$, and the
cell-periodic functions obey
\begin{equation}
	  \imath\hbar\,\partial_t u_{n\vb k}
	    = \left\{\frac{1}{2m_e}\Big[-\imath\hbar\nabla + \hbar\vb k
	          - \frac{q_e}{c}\vb A(t)\Big]^2
		      + V_{\mathrm{loc}}[\rho] + \hat V^{V}_{\mathrm{NL},\vb k}(t)\right\}u_{n\vb k},
		        \label{eq:VG_bloch}
\end{equation}
with
$\hat V^{V}_{\mathrm{NL},\vb k}
 = e^{-\imath\vb k\cdot\vb r}\,\hat V^{V}_{\mathrm{NL}}\,e^{\imath\vb k\cdot\vb r}$.
 Both the kinetic and the nonlocal terms depend on $\vb k$ and $\vb A(t)$ only
 through the shifted wave vector
 \begin{equation}
	   \vb k(t) = \vb k - \frac{q_e}{\hbar c}\,\vb A(t),
	     \qquad
	       \hat H^{V}_{\vb k}(t) = \hat H_{\vb k(t)}[\rho],
	         \label{eq:shifted_k}
 \end{equation}
 where $\hat H_{\vb k}$ is the field-free Bloch Hamiltonian. For the nonlocal
 term this follows from $e^{-\imath\vb k\cdot\vb r}\hat U_\chi
 = e^{-\imath\vb k(t)\cdot\vb r}$, so that
 $\hat V^{V}_{\mathrm{NL},\vb k} = \hat V_{\mathrm{NL},\vb k(t)}$. The homogeneous
 field thus rigidly shifts the crystal momentum of every Bloch state; for
 electrons, $\vb k(t) = \vb k + e\vb A(t)/(\hbar c)$. This permits propagation
 at fixed $\vb k$, with the nonlocal projectors evaluated at the shifted wave
 vector. Because the gauge phase is not lattice periodic,
 $\hat U_\chi(\vb r+\vb R,t)
  = e^{\imath q_e\vb A(t)\cdot\vb R/(\hbar c)}\,\hat U_\chi(\vb r,t)$,
  the transformation cannot be carried out within a periodic cell. It is used
  only to derive the VG Hamiltonian, after which the propagation is carried out
  entirely in the Bloch representation of Eq.~\eqref{eq:VG_bloch}.

%%%%%%%%%%%%%%%%%%%%%%%%%%%%%%%%%%%%%%%%%%%%%%%%%%%%

